\documentclass[10pt,a4paper]{amsart}
\usepackage{amsmath,amssymb,amsthm,graphicx}
% Retained preamble for the Stacks Project chapter snapshot in this directory.
% The chapter's own preamble.tex is a copy of the shared Stacks Project preamble; the block
% below is sliced MECHANICALLY, line for line, from it. Nothing was added, removed or
% redefined; the wrapper class, the xr-hyper/hyperref wiring and the externaldocument list
% are omitted because the assembler supplies the document class.
%
% lines taken from preamble.tex: 3, 5, 9, 11, 12, 156, 157, 158, ... (51 lines total)
%
% package coverage restored from the chapter's own preamble (2026-10-13): xy, xy-2cell

\usepackage[all]{xy}
\xyoption{2cell}
\UseAllTwocells

\newenvironment{reference}{\comment}{\endcomment}
\newenvironment{slogan}{\comment}{\endcomment}
\newenvironment{history}{\comment}{\endcomment}
\theoremstyle{plain}
\newtheorem{theorem}[subsection]{Theorem}
\newtheorem{proposition}[subsection]{Proposition}
\newtheorem{lemma}[subsection]{Lemma}
\theoremstyle{definition}
\newtheorem{definition}[subsection]{Definition}
\newtheorem{example}[subsection]{Example}
\newtheorem{exercise}[subsection]{Exercise}
\newtheorem{situation}[subsection]{Situation}
\theoremstyle{remark}
\newtheorem{remark}[subsection]{Remark}
\newtheorem{remarks}[subsection]{Remarks}
\numberwithin{equation}{subsection}
\def\lim{\mathop{\mathrm{lim}}\nolimits}
\def\colim{\mathop{\mathrm{colim}}\nolimits}
\def\Spec{\mathop{\mathrm{Spec}}}
\def\Hom{\mathop{\mathrm{Hom}}\nolimits}
\def\Ext{\mathop{\mathrm{Ext}}\nolimits}
\def\SheafHom{\mathop{\mathcal{H}\!\mathit{om}}\nolimits}
\def\SheafExt{\mathop{\mathcal{E}\!\mathit{xt}}\nolimits}
\def\Sch{\mathit{Sch}}
\def\Mor{\mathop{\mathrm{Mor}}\nolimits}
\def\Ob{\mathop{\mathrm{Ob}}\nolimits}
\def\Sh{\mathop{\mathit{Sh}}\nolimits}
\def\NL{\mathop{N\!L}\nolimits}
\def\CH{\mathop{\mathrm{CH}}\nolimits}
\def\proetale{{pro\text{-}\acute{e}tale}}
\def\etale{{\acute{e}tale}}
\def\QCoh{\mathit{QCoh}}
\def\Ker{\mathop{\mathrm{Ker}}}
\def\Im{\mathop{\mathrm{Im}}}
\def\Coker{\mathop{\mathrm{Coker}}}
\def\Coim{\mathop{\mathrm{Coim}}}
\DeclareMathSymbol{\boxtimes}{\mathbin}{AMSa}{"02}
\def\QCohstack{\mathcal{QC}\!\mathit{oh}}
\def\Cohstack{\mathcal{C}\!\mathit{oh}}
\def\Spacesstack{\mathcal{S}\!\mathit{paces}}
\def\Quotfunctor{\mathrm{Quot}}
\def\Hilbfunctor{\mathrm{Hilb}}
\def\Curvesstack{\mathcal{C}\!\mathit{urves}}
\def\Polarizedstack{\mathcal{P}\!\mathit{olarized}}
\def\Complexesstack{\mathcal{C}\!\mathit{omplexes}}
\def\Pic{\mathop{\mathrm{Pic}}\nolimits}
\def\Picardstack{\mathcal{P}\!\mathit{ic}}
\def\Picardfunctor{\mathrm{Pic}}
\def\Deformationcategory{\mathcal{D}\!\mathit{ef}}
\usepackage{cleveref}
\crefname{originallemma}{original lemma}{original lemmas}
\Crefname{originallemma}{Original lemma}{Original lemmas}
\makeatletter
\newlabel{homology-lemma-double-complex-gives-resolution}{{0133}{source}{}{sourcepointer-1}{}}
\newlabel{homology-lemma-double-complex-gives-resolution@cref}{{[originallemma][0][]0133}{[1][0][]source}{}{}{}}
\makeatother
\crefname{originallemma}{original lemma}{original lemmas}
\Crefname{originallemma}{Original lemma}{Original lemmas}
\makeatletter
\newlabel{homology-lemma-ss-double-complex}{{0130}{source}{}{sourcepointer-2}{}}
\newlabel{homology-lemma-ss-double-complex@cref}{{[originallemma][0][]0130}{[1][0][]source}{}{}{}}
\makeatother
\crefname{originallemma}{original lemma}{original lemmas}
\Crefname{originallemma}{Original lemma}{Original lemmas}
\makeatletter
\newlabel{homology-section-double-complex}{{012X}{source}{}{sourcepointer-3}{}}
\newlabel{homology-section-double-complex@cref}{{[originallemma][0][]012X}{[1][0][]source}{}{}{}}
\makeatother
\crefname{originallemma}{original lemma}{original lemmas}
\Crefname{originallemma}{Original lemma}{Original lemmas}
\makeatletter
\newlabel{lemma-injective-trivial-cech-variant}{{09VW}{source}{}{sourcepointer-4}{}}
\newlabel{lemma-injective-trivial-cech-variant@cref}{{[originallemma][0][]09VW}{[1][0][]source}{}{}{}}
\makeatother
\title{Hypercovering spectral sequence over a presheaf base}
\author{The Stacks Project Authors}
\date{stacks-project:hypercovering.tex:local-snapshot-20261012; GNU FDL 1.2 or later}
\begin{document}
\maketitle
\noindent\textit{Editorial note.} Statement span is lines 1398-1421 of sources/stacks-hypercovering/pinned-original-source.tex; proof span is lines 1423-1475, the complete author proof (51 body lines).\par\medskip
\noindent\textit{Editorial note.} Navigation stubs: 4 cited-but-not-included labels are supplied as non-load-bearing pointers with their pinned tags where the snapshot tag index has them.\par\medskip
\noindent\textit{Editorial note.} extra\_packages is empty on purpose: the retained preamble already loads xy with options and a second declaration clashes in the assembler wrapper.\par\medskip
\noindent\textit{Editorial note.} Text is retained verbatim, including the author's own typos and placeholders.\par\medskip
\noindent\textit{Editorial note.} Queue index 7 of queue\_hypercovering.json (pb=2243, stmt 1398-1421, proof 1423-1475).\par\medskip
\noindent\textit{Editorial note.} stacks-hypercovering wave-1 (hc01): fresh adjudication accept (2026-10-14); pool per standing fallback order (hypercovering -> pic -> spaces-more-cohomology) after md01 closed the moduli pool; hypercovering.tex 53 labels/34 with proof judged 26 accept / 8 reject / 14 pb=0 env-dup rows excluded; spans cite-free and reference/slogan/history-free so no patchenv needed.\par\medskip
\noindent\textit{Allowed standard prerequisites (editorial).}\begin{itemize}
\item Explicit prerequisite boundary. The cited inputs are prerequisites; the argument in the retained span is not replaced by any cited result.
\end{itemize}\medskip
\ifdefined\hypertarget\hypertarget{sourcepointer-1}{}\fi
\noindent\textit{Source reference.} Original lemma 0133: sources/stacks-hypercovering/pinned-original-source.tex, source lines 1475--1475. Cited by the retained statement or proof but not part of the collected spans; supplied only as a navigation pointer. The argument does not depend on its content.\par\medskip
\ifdefined\hypertarget\hypertarget{sourcepointer-2}{}\fi
\noindent\textit{Source reference.} Original lemma 0130: sources/stacks-hypercovering/pinned-original-source.tex, source lines 1475--1475. Cited by the retained statement or proof but not part of the collected spans; supplied only as a navigation pointer. The argument does not depend on its content.\par\medskip
\ifdefined\hypertarget\hypertarget{sourcepointer-3}{}\fi
\noindent\textit{Source reference.} Original lemma 012X: sources/stacks-hypercovering/pinned-original-source.tex, source lines 1475--1475. Cited by the retained statement or proof but not part of the collected spans; supplied only as a navigation pointer. The argument does not depend on its content.\par\medskip
\ifdefined\hypertarget\hypertarget{sourcepointer-4}{}\fi
\noindent\textit{Source reference.} Original lemma 09VW: sources/stacks-hypercovering/pinned-original-source.tex, source lines 1364--1376. Cited by the retained statement or proof but not part of the collected spans; supplied only as a navigation pointer. The argument does not depend on its content.\par\medskip
\begin{lemma}
\label{lemma-cech-spectral-sequence-variant}
Let $\mathcal{C}$ be a site with equalizers and fibre products.
Let $\mathcal{G}$ be a presheaf on $\mathcal{C}$.
Let $K$ be a hypercovering of $\mathcal{G}$.
Let $\mathcal{F}$ be a sheaf of abelian groups on $\mathcal{C}$.
There is a map
$$
s(\mathcal{F}(K)) \longrightarrow R\Gamma(\mathcal{G}, \mathcal{F})
$$
in $D^{+}(\textit{Ab})$ functorial in $\mathcal{F}$, which induces
a natural transformation
$$
\check{H}^i(K, -) \longrightarrow H^i(\mathcal{G}, -)
$$
of functors $\textit{Ab}(\mathcal{C}) \to \textit{Ab}$. Moreover,
there is a spectral sequence $(E_r, d_r)_{r \geq 0}$ with
$$
E_2^{p, q} = \check{H}^p(K, \underline{H}^q(\mathcal{F}))
$$
converging to $H^{p + q}(\mathcal{G}, \mathcal{F})$.
This spectral sequence is functorial in $\mathcal{F}$ and
in the hypercovering $K$.
\end{lemma}

\begin{proof}
Choose an injective resolution $\mathcal{F} \to \mathcal{I}^\bullet$
in the category of abelian sheaves on $\mathcal{C}$. Consider the
double complex $A^{\bullet, \bullet}$ with terms
$$
A^{p, q} = \mathcal{I}^q(K_p)
$$
where the differential $d_1^{p, q} : A^{p, q} \to A^{p + 1, q}$
is the one coming from the differential $\mathcal{I}^p \to \mathcal{I}^{p + 1}$
and the differential $d_2^{p, q} : A^{p, q} \to A^{p, q + 1}$ is the
one coming from the differential on the complex
$s(\mathcal{I}^p(K))$ associated to the cosimplicial abelian group
$\mathcal{I}^p(K)$ as explained above.
We will use the two spectral
sequences $({}'E_r, {}'d_r)$ and $({}''E_r, {}''d_r)$
associated to this double complex, see
Homology, Section \ref{homology-section-double-complex}.

\medskip\noindent
By Lemma \ref{lemma-injective-trivial-cech-variant} the complexes
$s(\mathcal{I}^p(K))$ are acyclic in positive degrees and have
$H^0$ equal to $H^0(\mathcal{G}, \mathcal{I}^p)$. Hence by
Homology, Lemma \ref{homology-lemma-double-complex-gives-resolution}
and its proof the spectral sequence $({}'E_r, {}'d_r)$ degenerates,
and the natural map
$$
H^0(\mathcal{G}, \mathcal{I}^\bullet) \longrightarrow
\text{Tot}(A^{\bullet, \bullet})
$$
is a quasi-isomorphism of complexes of abelian groups. The map
$s(\mathcal{F}(K)) \longrightarrow R\Gamma(\mathcal{G}, \mathcal{F})$
of the lemma is the composition of the natural map
$s(\mathcal{F}(K)) \to \text{Tot}(A^{\bullet, \bullet})$
followed by the inverse of the displayed quasi-isomorphism above.
This works because $H^0(\mathcal{G}, \mathcal{I}^\bullet)$
is a representative of $R\Gamma(\mathcal{G}, \mathcal{F})$.

\medskip\noindent
Consider the spectral sequence $({}''E_r, {}''d_r)_{r \geq 0}$. By
Homology, Lemma \ref{homology-lemma-ss-double-complex}
we see that
$$
{}''E_2^{p, q} = H^p_{II}(H^q_I(A^{\bullet, \bullet}))
$$
In other words, we first take cohomology with respect to
$d_1$ which gives the groups
${}''E_1^{p, q} = \underline{H}^p(\mathcal{F})(K_q)$.
Hence it is indeed the case (by the description of the differential
${}''d_1$) that
${}''E_2^{p, q} = \check{H}^p(K, \underline{H}^q(\mathcal{F}))$.
Since this spectral sequence converges to the cohomology of
$\text{Tot}(A^{\bullet, \bullet})$ the proof is finished.
\end{proof}

\end{document}
